\section{Research Directions 与 Q\&A：仍不知道什么}

课堂最后没有用“scale solves everything”收尾，而是列出 uncertainty、inference-time learning、adaptive thinking、planning 和 quantization。Q\&A 又补上 non-autoregressive ordering、language world knowledge、generalization、multi-agent coordination、modularity 与 gradient descent 等限制条件。这些 spoken nodes 是本讲 teacher voice 最密集的部分。

\subsection{五条公开 Research Directions}

第一，模型应知道自己不知道什么；第二，模型能否在 inference time 学会新 skill，而不是每次都更新参数；第三，能否按问题难度分配 adaptive compute；第四，reasoning/planning 是否可由小模型、external planner 或更好 representation 完成；第五，低 precision 能否继续降低 memory traffic 与提高 matrix throughput。

\lecturefigure{slide-43-research-directions.jpg}{不确定性、inference-time learning、adaptive thinking、planning 与 quantization}{官方录像恢复幻灯片 V043；01:00:51}

\begin{knowledgebox}{读图：五个方向跨越 model、algorithm 与 system}
Knowing what you do not know 属于 calibration/evaluation；learning skills at inference time 属于 memory/adaptation；adaptive thinking 和 planning 属于 compute allocation；quantization 属于 numerical/system co-design。它们不能被一个“更大模型”指标统一验收。
\end{knowledgebox}

\teachervoice{讲者注意到 diffusion model 可通过更多采样 compute 改善质量，而语言模型当时缺少同样清晰的 adaptive-compute knob。他追问：能否让小模型在难题上多想、连接 external planner，从而接近大模型 reasoning，而不是每个 token 都支付最大成本。}

\subsection{Non-autoregressive Decoding：真正缺的是 Order Oracle}

Q\&A 进一步解释 parallel generation。若已经生成的 latent/context 足以使下一组 token 条件独立，就可并行产生这一组；难点是学习这个 grouping/order。一个分组 factorization 为：
\[
p(y\mid x)=\prod_{g=1}^{G}\prod_{i\in S_g}p\!\left(y_i\mid x,y_{S_{<g}}\right),
\]
其中 $S_g$ 是第 $g$ 个可并行生成的 token group，$S_{<g}$ 是之前已生成 groups。若 oracle 给出正确 groups，问题容易许多；实际模型还要发现 conditional independence，并处理多模态输出分布。

\teachervoice{讲者说如果有人给出“先生成这三个词，再生成那两个词”的 oracle ordering，non-autoregressive generation 会更有机会。他们尝试用 vector-quantized latent sequence 建模依赖，translation 上可工作，却需要 distillation 降低数据 entropy，最终 practical impact 不及 speculative decoding。}

\subsection{Language-only World Model：不充分，但比直觉更有用}

课堂没有接受“文本能学到全部物理世界”，也没有否认其 usefulness。大语言模型从文本中吸收大量 object、procedure 与 social knowledge，可作为 robotics planner 的高层先验；perception、control 与 environment feedback 仍由其他系统提供。关键是区分 knowledge coverage 与 action grounding。

\begin{warningbox}{能规划不等于能闭环控制}
语言模型说出合理步骤，只证明它能组织文本先验；机器人还需识别当前对象、估计状态、满足动力学、安全执行并从失败恢复。World knowledge 是 planner input，不是 physical verification。
\end{warningbox}

\teachervoice{Vaswani 的立场很务实：language alone 也许不完整，但机器人已能把大模型当 planner，说明文本承载的世界信息比许多人预期更多。我们甚至还没有完全榨取现有模型的 usefulness，因此不能用终极不足否定当前工程价值。}

\subsection{Generalization：新能力可能来自大规模 Recombination}

如果训练数据覆盖极广，传统 OOD 边界变得难定义。模型可能没有见过“用 Chaucer 韵律描述这堂 Stanford lecture”的精确样本，却见过讲座、Chaucer、韵律和改写，能够组合成新输出。这是有用的 compositional generalization，即使其原子信息来自训练分布。

\begin{knowledgebox}{不要把 Recombination 贬成“只是记忆”}
检验应区分 exact memorization、template interpolation、systematic composition 与 algorithmic extrapolation。大规模 recombination 可以产生真实价值，却不保证长度、数字、规则或状态空间上的任意外推；不同 benchmark 测量不同 generalization。
\end{knowledgebox}

\subsection{Multi-agent Systems：Decomposition、Coordination、Verification}

多个 agent 是否优于一个大模型，不由角色数量决定。讲者给出三个基础问题：目标如何按各 agent skill 分解，过程如何协调共享状态和依赖，结果如何验证。没有 verification，更多 agent 只会制造更多文本与更复杂 failure propagation。

\begin{importantbox}{Multi-agent 的三项最低验收}
\textbf{Goal decomposition} 要把任务拆成可判定 subgoal；\textbf{coordination} 要处理顺序、资源、冲突与共享 memory；\textbf{verification} 要用测试、环境状态或外部证据检查结果。三者缺一，所谓 swarm 只是并行采样。
\end{importantbox}

\teachervoice{讲者半开玩笑说“最好的 agents 是 neurons”，因为它们通过 gradient 协同更新。对于显式 multi-agent，他没有给出万能架构，而是明确承认 decomposition、coordination、verification 仍处早期。}

\subsection{Modularity、MoE 与“Architecture 服务 Gradient Descent”}

人类设计的 module 常因 routing、credit assignment 和联合优化困难而失败。Mixture of Experts（MoE）提供 learned routing，使部分 specialization 在 end-to-end gradient descent 中出现；但 expert identity 不等于完整认知模块，责任仍可能分布在 attention、residual stream 与多个 experts 中。

\teachervoice{Vaswani 说自己有时觉得“我们只是在构建 architecture 来服务 gradient descent”。这句话不是贬低架构，而是指出可训练性是硬约束：一个模块划分再优雅，如果 credit 无法传播、routing 无法学习，就难以替代较均质但可优化的网络。}

\subsection{本章小结}

Research directions 跨越 calibration、adaptation、compute allocation、planning 和 numerical efficiency。Q\&A 进一步表明：parallel generation 需要 order/independence structure，language world knowledge 有用但不等于 grounding，multi-agent 需要 verification，modularity 必须服从可训练性。

